\chapter{Εισαγωγή}
\label{ch:intro}

Στο κεφάλαιο αυτό παρουσιάζεται το πλαίσιο που καθιστά αναγκαία την πρόβλεψη τιμών και
φορτίου στις ηλεκτρικές αγορές, περιγράφεται το ιδιαίτερο πλαίσιο της ελληνικής αγοράς,
διατυπώνονται τα ερευνητικά ερωτήματα που καθοδηγούν την εργασία και παρουσιάζεται
η δομή των κεφαλαίων που ακολουθούν.

% ─────────────────────────────────────────────────────────────────────────────
\section{Ενεργειακή Μετάβαση και Ηλεκτρικές Αγορές}
\label{sec:energy_transition}
% ─────────────────────────────────────────────────────────────────────────────

Η παγκόσμια ενεργειακή μετάβαση αποτελεί μια από τις βαθύτερες αναδιαρθρώσεις του
ενεργειακού συστήματος που έχουν καταγραφεί ιστορικά. Η κλιματική αλλαγή, οι
δεσμευτικοί στόχοι μηδενικών εκπομπών άνθρακα για το 2050 και η ραγδαία μείωση του
κόστους φωτοβολταϊκών και αιολικών εγκαταστάσεων αναγκάζουν τα εθνικά δίκτυα να
μεταβούν από ένα κεντρικό μοντέλο ορυκτών καυσίμων σε ένα αποκεντρωμένο,
ευέλικτο και χαμηλών εκπομπών μοντέλο~\cite{weron2014electricity}.

Παράλληλα, η απελευθέρωση των ηλεκτρικών αγορών — που ξεκίνησε στα τέλη του 20ού
αιώνα με πρωτοπόρες τις αγγλοσαξονικές χώρες — επέτρεψε την εισαγωγή του
ανταγωνισμού σε κλάδους που ιστορικά λειτουργούσαν ως φυσικά μονοπώλια. Σήμερα,
στην Ευρώπη, η χονδρεμπορική αγορά ηλεκτρικής ενέργειας οργανώνεται κυρίως γύρω
από την \textit{Αγορά Επόμενης Ημέρας} (\gls{dam}), όπου αντιστοιχίζονται προσφορές
παραγωγής και ζήτησης για κάθε ωριαία περίοδο της επόμενης
ημέρας~\cite{cacm2015}. Σε αυτές τις αγορές, η αποτελεσματική κατανομή πόρων
βασίζεται ολοένα και περισσότερο σε ακριβείς βραχυπρόθεσμες προβλέψεις.

\subsection*{Πρόβλεψη Τιμής Ηλεκτρικής Ενέργειας}

Η \textit{Πρόβλεψη Τιμής Ηλεκτρικής Ενέργειας} (\gls{epf}) αναδύεται ως κρίσιμο
εργαλείο για τη διαχείριση ρίσκου και τη βελτιστοποίηση θέσης στις ηλεκτρικές
αγορές~\cite{lago2021}. Η τιμή εκκαθάρισης της αγοράς επόμενης ημέρας (DAM
clearing price) εξαρτάται από τη συνάρτηση οριακού κόστους παραγωγής, την παρεμβολή
ΑΠΕ μηδενικού οριακού κόστους, τις τιμές καυσίμων και CO\textsubscript{2}, τη
ζήτηση φορτίου και τη δυναμικότητα διασυνδέσεων. Η συνδυαστική δράση αυτών των
παραγόντων καθιστά την τιμή ηλεκτρισμού ιδιαίτερα \textit{μη-γραμμική},
\textit{εποχιακή} και \textit{ευμετάβλητη}, με σποραδικές τιμές-ακραίες τιμές
(spikes) που μπορούν να υπερβαίνουν κατά 10--20× τη μέση τιμή.

Η ακριβής πρόβλεψη παρέχει ανταγωνιστικό πλεονέκτημα σε εμπόρους ενέργειας,
παραγωγούς, βιομηχανικούς καταναλωτές και συσσωρευτές (aggregators) για τη βέλτιστη
ανάθεση των προσφορών τους στην DAM αλλά και στην Αγορά Εξισορρόπησης.

\subsection*{Πρόβλεψη Ηλεκτρικού Φορτίου}

Η \textit{Πρόβλεψη Ηλεκτρικού Φορτίου} (\gls{elf}) προηγείται χρονολογικά της EPF
ως αντικείμενο ακαδημαϊκής έρευνας~\cite{hong2020}. Η ακριβής εκτίμηση της
ζήτησης αποτελεί τον πυρήνα του προγραμματισμού παραγωγής για τους Διαχειριστές
Συστήματος Μεταφοράς (\gls{tso}), ελαχιστοποιώντας το κόστος εξισορρόπησης και
τους κινδύνους ακάλυπτης ζήτησης. Η ηλεκτρική κατανάλωση χαρακτηρίζεται από
σαφή ημερήσια και εβδομαδιαία εποχικότητα (πρωινές και βραδινές αιχμές, χαμηλά
φορτία Κυριακής), ενώ η επίδραση της θερμοκρασίας είναι κυρίαρχη.

Εκτός από τη χρήση του φορτίου ως αυτόνομης πρόβλεψης, η ημερήσια πρόγνωση
ηλεκτρικής ζήτησης (\textit{day-ahead load forecast}) χρησιμεύει και ως \textit{είσοδος
δύο σταδίων} στην EPF: οι προβλέψεις φορτίου δημοσιεύονται πριν την κλεισίδωση
της DAM και ενσωματώνονται ως χαρακτηριστικά στα μοντέλα τιμής.


% ─────────────────────────────────────────────────────────────────────────────
\section{Η Ελληνική Αγορά Ηλεκτρισμού}
\label{sec:greek_market}
% ─────────────────────────────────────────────────────────────────────────────

Η ελληνική αγορά ηλεκτρισμού εντάχθηκε στο ευρωπαϊκό μοντέλο Target Model την
1η Νοεμβρίου 2020, εναρμονιζόμενη με τον κανονισμό CACM~\cite{cacm2015} και
υιοθετώντας τον αλγόριθμο EUPHEMIA για τη συνδυαστική εκκαθάριση τιμών στο
πλαίσιο του \gls{sdac}~\cite{nemo_sdac}. Η γεωγραφική θέση της Ελλάδας στο
νοτιοανατολικό άκρο της ευρωπαϊκής ηπείρου, με διασυνδέσεις προς Βουλγαρία,
Ιταλία, Αλβανία, Βόρεια Μακεδονία και εσωτερικές νησιωτικές διασυνδέσεις,
προσδίδει ιδιαίτερο χαρακτήρα στη διαμόρφωση της τιμής.

Ο \gls{admhe} ασκεί τον ρόλο του TSO, εποπτεύοντας τη λειτουργία της Αγοράς
Εξισορρόπησης και διαχειριζόμενος την ασφαλή λειτουργία του Ελληνικού
Διασυνδεδεμένου Συστήματος~\cite{ipto_balancing_market}. Παράλληλα, ο ΔΑΠΕΕΠ
διαχειρίζεται τη δραστηριότητα προμήθειας και εκκαθάρισης αγοράς.

\subsection*{Ενεργειακό Μίγμα και Τιμές}

Το ενεργειακό μίγμα της Ελλάδας χαρακτηρίζεται από αυξανόμενη συμμετοχή
\gls{res}: η ηλιακή ισχύς υπερέβη τα 8~GW εγκατεστημένης ισχύος στα τέλη του
2024, ενώ η αιολική συνεισφορά ανέρχεται σε ~5~GW. Ταυτόχρονα, το φυσικό αέριο
αποτελεί τον βασικό ρυθμιστή (marginal fuel) στο 30--50\% των ωρών~\cite{ypen_target_model_2020}.
Αυτός ο συνδυασμός δημιουργεί μια αγορά όπου η τιμή επηρεάζεται ταυτόχρονα από:
\begin{itemize}
    \item τη διεθνή τιμή φυσικού αερίου (TTF/NBP),
    \item τη διεθνή τιμή δικαιωμάτων εκπομπής CO\textsubscript{2} (EUA),
    \item την ηλιακή και αιολική παραγωγή (που εισάγει αβεβαιότητα),
    \item τις ηλεκτρικές διασυνδέσεις με γειτονικές αγορές.
\end{itemize}

Κατά την περίοδο αξιολόγησης (Δεκέμβριος 2025 -- Φεβρουάριος 2026), οι τιμές
κυμάνθηκαν σε εύρος 30--200 €/MWh, με μέση τιμή~110,4~€/MWh (Δεκέμβριος 2025),
ενώ τον Φεβρουάριο 2026 παρατηρήθηκε έντονη πτώση της μέσης τιμής στα~77,8~€/MWh
(αλλαγή καθεστώτος τιμών) με ταυτόχρονη αύξηση της τυπικής
απόκλισης (33,6 $\to$ 45,4 €/MWh).

\begin{figure}[htbp]
    \centering
    \includegraphics[width=\textwidth]{ch6_forecast_analysis/01_ts_price.png}
    \caption{Χρονοσειρά τιμής DAM (€/MWh) για την Ελλάδα — Δεκέμβριος 2025 έως
    Φεβρουάριος 2026. Οι σκιαγραφημένες περιοχές δείχνουν τον Ιανουάριο και
    Φεβρουάριο 2026, στους οποίους παρατηρείται αυξημένη μεταβλητότητα και πτώση
    του επιπέδου τιμής.}
    \label{fig:ts_price}
\end{figure}

\begin{figure}[htbp]
    \centering
    \includegraphics[width=\textwidth]{ch6_forecast_analysis/02_ts_load.png}
    \caption{Χρονοσειρά φορτίου (MW) για την Ελλάδα — Δεκέμβριος 2025 έως
    Φεβρουάριος 2026. Το χειμερινό φορτίο εμφανίζει σαφή ημερήσια εποχικότητα
    και σταδιακή μείωση καθώς πλησιάζει η άνοιξη.}
    \label{fig:ts_load}
\end{figure}

\subsection*{Βιβλιογραφικό Πλαίσιο για την Ελλάδα}

Η επιστημονική βιβλιογραφία για την EPF/ELF στην ελληνική αγορά είναι περιορισμένη
σε σχέση με τις αναπτυγμένες αγορές (Γερμανία, Αυστραλία, Ισπανία). Ο
Κανούσης~\cite{kanousis2025} εφαρμόζει XGBoost/LGBM στην ελληνική DAM τιμή για
την περίοδο 2018--2023, αποδεικνύοντας την ανωτερότητα των boosting μεθόδων έναντι
ARIMA/SVR. Ο Μαγκλαράς~\cite{magklaras2022} επικεντρώνεται στη βραχυπρόθεσμη
πρόβλεψη φορτίου με ανάλυση εποχικότητας και μοντέλα βαθιάς μάθησης.
Η παρούσα εργασία επεκτείνει αυτά τα αποτελέσματα καλύπτοντας και τις δύο
εργασίες, συγκρίνοντας τέσσερις στρατηγικές πρόβλεψης και αξιολογώντας σε
πολλαπλές χρονικές περιόδους με ακτυαλιστικά δεδομένα 2025--2026.


% ─────────────────────────────────────────────────────────────────────────────
\section{Κίνητρο και Ερευνητικά Ερωτήματα}
\label{sec:motivation}
% ─────────────────────────────────────────────────────────────────────────────

Παρά τη σημαντική πρόοδο στην EPF/ELF, διάφορα ερευνητικά κενά παραμένουν
ανοικτά. Πρώτον, οι περισσότερες εργασίες αξιολογούν ένα ή δύο μοντέλα σε μία
μόνο στρατηγική πρόβλεψης (συνήθως teacher-forcing), παραβλέποντας τη σύγκριση
μεταξύ TF, Recursive, MIMO και Walk-Forward Retraining. Δεύτερον, σπάνια
αξιολογείται ταυτόχρονα η πρόβλεψη τιμής \textit{και} φορτίου υπό κοινό
πειραματικό πλαίσιο. Τρίτον, η επίδοση σε μεγάλες χρονικές περιόδους (3 μήνες)
με πιθανή \textit{μεταβολή κατανομής} (distributional shift) παραμένει
αναξιολόγητη σε πολλές εργασίες.

\medskip
Οι κύριες ερευνητικές ερωτήσεις που καθοδηγούν την παρούσα εργασία είναι:

\begin{enumerate}[label=\textbf{ΕΕ\arabic*.}]
    \item Ποια στρατηγική πρόβλεψης (Teacher-Forcing, Recursive, MIMO, Monthly
    Retraining) επιτυγχάνει τη μικρότερη MAE σε επίπεδο τιμής DAM και ηλεκτρικού
    φορτίου στην ελληνική αγορά;

    \item Πώς επηρεάζει ο ορίζοντας αξιολόγησης (1 εβδομάδα, 1 μήνας, 3 μήνες)
    την κατάταξη των στρατηγικών και αλγόριθμων;

    \item Υπάρχουν συστηματικές στατιστικά σημαντικές διαφορές επίδοσης μεταξύ των
    μοντέλων, ή παρατηρούνται αλληλεπικαλύψεις λόγω μεταβλητότητας του σήματος;

    \item Ποια χαρακτηριστικά (lags, εποχικότητα, εξωγενείς παράμετροι) αξιοποιούν
    περισσότερο τα μοντέλα ανά στρατηγική;

    \item Μπορεί η μηνιαία επανεκπαίδευση (walk-forward retraining) να αντισταθμίσει
    τη μεταβολή κατανομής κατά τη διάρκεια τεταρτημορίων υψηλής μεταβλητότητας
    (Q1 2026);
\end{enumerate}

\medskip
Η διατύπωση αυτών των ερωτημάτων οδήγησε στη σχεδίαση ενός ολοκληρωμένου
πειραματικού πλαισίου με πέντε αλγόριθμους ML, τέσσερις στρατηγικές πρόβλεψης,
τρεις χρονικές περιόδους αξιολόγησης και στατιστική επαλήθευση DM test.


% ─────────────────────────────────────────────────────────────────────────────
\section{Σκοπός και Κύριες Συνεισφορές}
\label{sec:purpose}
% ─────────────────────────────────────────────────────────────────────────────

Ο σκοπός της εργασίας είναι η \textit{συστηματική, συγκριτική αξιολόγηση
στρατηγικών βραχυπρόθεσμης πρόβλεψης} τιμής DAM και ηλεκτρικού φορτίου στην
ελληνική αγορά, με έμφαση στη μεθοδολογική αυστηρότητα: αποφυγή διαρροής
πληροφορίας (data leakage), ακέραια αξιολόγηση εκτός δείγματος (out-of-sample)
και στατιστικά έγκυρη σύγκριση.

\subsection*{Κύριες Συνεισφορές}

\begin{enumerate}
    \item \textbf{Σύγκριση τεσσάρων στρατηγικών πρόβλεψης} σε κοινό πειραματικό
    πλαίσιο: Teacher-Forcing (TF), Recursive/Open-Loop (OL) με Scheduled Sampling,
    Multi-Input Multi-Output (MIMO) με χρήση dense lags, και Walk-Forward
    Monthly Retraining. Κάθε στρατηγική εφαρμόζεται σε πέντε αλγόριθμους ΜΜ
    (LightGBM, XGBoost, Random Forest, SVR, MLP) για δύο εργασίες.

    \item \textbf{Πολυεπίπεδη αξιολόγηση} σε τρεις χρονικές κλίμακες (Δεκέμβριος
    2025 = 744h, Q4 2025 = 2208h, Q1 2026 = 2160h), αποκαλύπτοντας πώς η σχετική
    κατάταξη στρατηγικών και αλγόριθμων μεταβάλλεται με τον ορίζοντα και τη
    μεταβλητότητα της αγοράς.

    \item \textbf{Walk-Forward Monthly Retraining} (μηνιαία επανεκπαίδευση):
    αποδεικνύεται ότι η επανεκπαίδευση κάθε μήνα μειώνει το \gls{mae} τιμής
    κατά 27,7\% (14.478~€/MWh $\to$ 10.463~€/MWh) και
    του φορτίου κατά 30,0\% (96,8~MW $\to$ 67,7~MW) σε σχέση με στατική εκπαίδευση
    κατά την Q1 2026 — η σημαντικότερη βελτίωση που τεκμηριώνεται στην εργασία.

    \item \textbf{Dense Lag Features για MIMO}: Η προσθήκη ενδοημερήσιων χρονικών
    υστερήσεων (lags 4--23) στις MIMO στρατηγικές βελτιώνει σημαντικά την
    ακρίβεια: \gls{mlp} −22,3\%, XGBoost −7,9\%, \gls{lgbm} −6,5\%, RF −2,3\%.
    Ταυτόχρονα, αποδεικνύεται ότι τα ίδια dense lags \textit{βλάπτουν} τη
    Recursive στρατηγική (+2,8\% για LGBM) λόγω συσσώρευσης σφάλματος.

    \item \textbf{Scheduled Sampling (SS)}: Αξιολογείται η τεχνική που εκπαιδεύει
    το μοντέλο σε μικτές εισόδους (πραγματικές τιμές + αυτοπροβλέψεις), αποδεικνύοντας
    βελτίωση έναντι plain Recursive: XGBoost τιμή −1,5\%, LGBM φορτίο −4,6\% vs
    plain LGBM OL.

    \item \textbf{Ανάλυση Σημαντικότητας Χαρακτηριστικών} (feature importance):
    Αναλύονται 15 παραλλαγές μοντέλων ανά αλγόριθμο και στρατηγική, αποκαλύπτοντας
    την μεταβολή της αξιοποίησης χαρακτηριστικών: υψηλή συγκέντρωση στο
    $y_{t-1}$ (88,3\%) για TF, ισότητα μεταξύ lag-24/48/168 για MIMO, ανάδυση
    αιολικής παραγωγής σε Direct-LGBM-Dense.

    \item \textbf{Στατιστικός Έλεγχος Diebold-Mariano} με διόρθωση
    Harvey-Leybourne-Newbold (HLN) για πλήρη πίνακα σύγκρισης 18 μοντέλων ανά
    εργασία (324 ζεύγη), ώστε να διακριθούν στατιστικά σημαντικές από
    τυχαίες διαφορές επίδοσης.

    \item \textbf{Διαδραστικό Dashboard}: Αναπτύσσεται αυτόνομο σύστημα
    οπτικοποίησης (standalone HTML/JavaScript) που επιτρέπει διαδραστική εξερεύνηση
    αποτελεσμάτων ανά στρατηγική, μοντέλο, εργασία και χρονική περίοδο, χωρίς
    απαίτηση web server ή προγραμματιστικών εργαλείων.
\end{enumerate}


% ─────────────────────────────────────────────────────────────────────────────
\section{Δομή της Εργασίας}
\label{sec:structure}
% ─────────────────────────────────────────────────────────────────────────────

Η εργασία οργανώνεται σε εννέα κεφάλαια ως εξής:

\begin{description}[leftmargin=0pt, style=nextline]
    \item[\textbf{Κεφάλαιο~\ref{ch:literature} — Βιβλιογραφική Ανασκόπηση:}]
    Παρουσιάζονται οι κυριότεροι αλγόριθμοι ΜΜ που χρησιμοποιούνται στην EPF/ELF,
    οι στρατηγικές πολυβήματης πρόβλεψης (TF, Recursive, MIMO, Walk-Forward), η
    τεχνική Scheduled Sampling, η βελτιστοποίηση υπερπαραμέτρων και το πλαίσιο της
    ελληνικής αγοράς στη βιβλιογραφία.

    \item[\textbf{Κεφάλαιο~\ref{ch:data} — Δεδομένα και Προεπεξεργασία:}]
    Περιγράφονται οι πηγές δεδομένων (ENTSO-E Transparency Platform), η διαδικασία
    συγχώνευσης και επεξεργασίας, η κατασκευή χαρακτηριστικών (lags, rolling
    statistics, ημερολογιακές μεταβλητές, εξωγενείς παράμετροι), και η πολιτική
    αποφυγής διαρροής πληροφορίας.

    \item[\textbf{Κεφάλαιο~\ref{ch:methods} — Μέθοδοι Πρόβλεψης:}]
    Παρουσιάζονται αναλυτικά οι αλγόριθμοι ML (LightGBM, XGBoost, RF, SVR, MLP),
    οι τέσσερις στρατηγικές πρόβλεψης με μαθηματικές εξισώσεις, η βελτιστοποίηση
    Optuna (OL-Aware, MIMO, CachedOptuna) και οι μετρικές αξιολόγησης (MAE, RMSE,
    sMAPE).

    \item[\textbf{Κεφάλαιο~\ref{ch:experiments} — Αποτελέσματα Αξιολόγησης:}]
    Παρουσιάζονται τα πειραματικά αποτελέσματα σε τρεις χρονικές περιόδους και
    τέσσερις στρατηγικές, με αναλυτικούς πίνακες επίδοσης, διαγράμματα σύγκρισης
    και ανάλυση των κυριότερων ευρημάτων (αποτυχία SVR, ανωτερότητα MR, dense lags
    paradox).

    \item[\textbf{Κεφάλαιο~\ref{ch:fi} — Σημαντικότητα Χαρακτηριστικών:}]
    Αναλύεται η σημαντικότητα χαρακτηριστικών ανά αλγόριθμο (gain/Gini-based) και
    στρατηγική, παρουσιάζεται η μεταβολή στην αξιοποίηση χαρακτηριστικών μεταξύ
    TF, Recursive και MIMO, και εξάγονται συμπεράσματα για τη σχεδίαση feature sets.

    \item[\textbf{Κεφάλαιο~\ref{ch:dm} — Στατιστικοί Έλεγχοι:}]
    Εφαρμόζεται ο έλεγχος Diebold-Mariano με HLN διόρθωση για πλήρη πίνακα
    σύγκρισης 18 μοντέλων, επιτρέποντας τη διάκριση στατιστικά σημαντικών
    διαφορών από τυχαίες μεταβολές.

    \item[\textbf{Κεφάλαιο~\ref{ch:dashboard} — Σύστημα Οπτικοποίησης:}]
    Περιγράφεται η αρχιτεκτονική και η λειτουργικότητα του διαδραστικού dashboard
    που αναπτύχθηκε για την οπτικοποίηση και εξερεύνηση των αποτελεσμάτων.

    \item[\textbf{Κεφάλαιο~\ref{ch:conclusions} — Συμπεράσματα:}]
    Συνοψίζονται τα κύρια ευρήματα της εργασίας, συζητούνται οι περιορισμοί και
    προτείνονται κατευθύνσεις για μελλοντική έρευνα.
\end{description}
